\section{Resource compilers: descent, windows, and complete lists}
\label{sec:resources}

A finite algorithm is useful to iPRA only when its stopping resource is
primitive recursive and its correctness can be maintained by a
quantifier-free invariant.  This section collects three generic resource
constructions that resolve recurring gaps in theorem packets.

\subsection{Fixed-dimensional lexicographic descent}

Let a state carry a tuple
\[
  a=(a_0,\ldots,a_{d-1}),\qquad 0\leq a_i\leq H_i(\vec x),
\]
where the dimension \(d\) is externally fixed.  Suppose every nonterminal step
strictly decreases the tuple in lexicographic order, while lower-priority
coordinates may reset within their bounds.

Define the mixed-radix rank
\[
 \rho(a)=a_0\prod_{j=1}^{d-1}(H_j+1)
 +a_1\prod_{j=2}^{d-1}(H_j+1)+\cdots+a_{d-1}.
\]

\begin{lemma}[Mixed-radix descent]
\label{lem:mixed-radix}
If \(b<_{\mathrm{lex}}a\) and both tuples satisfy the coordinate bounds, then
\(\rho(b)<\rho(a)\).  Consequently every strict lexicographic descent has
length less than
\[
  \prod_{i=0}^{d-1}(H_i+1).
\]
\end{lemma}

\begin{proof}
Let \(r\) be the first coordinate where \(a\) and \(b\) differ.  Then
\(b_r\leq a_r-1\).  Decreasing coordinate \(r\) lowers the rank by at least
\(\prod_{j>r}(H_j+1)\).  The largest possible increase in all later
coordinates is
\[
 \sum_{i>r}H_i\prod_{j>i}(H_j+1)
 =\prod_{j>r}(H_j+1)-1.
\]
The net change is therefore at most \(-1\).  The rank is nonnegative and below
the product of radices, so the chain-length bound follows.
\end{proof}

All products are primitive recursive for fixed \(d\).  The lemma yields a
fuel term for Euclidean descent, reduction algorithms, finite rewrite systems,
triangular elimination, and bounded searches with resetting counters.

If coordinate bounds depend monotonically on time or earlier states, replace
\(H_i\) by a proved envelope.  More sophisticated nested-iteration bounds from
the resources packet can be used when a simple product is too weak.  The
fixed-dimension condition must remain visible: arbitrary input-dependent
nesting depth can escape primitive recursion.

\subsection{Simultaneous finite-window selection}

Suppose a finite interval \([0,L)\) has exceptional predicates
\(E_1,\ldots,E_m\).  Assume each has a proved counting budget
\[
 \#\{i<L:E_j(i)\}\leq b_j.
\]

\begin{theorem}[Finite union window]
\label{thm:finite-window}
If \(L>\sum_{j=1}^m b_j\), bounded least search constructs an
\(i<L\) satisfying none of the \(E_j\).  iPRA proves the witness's goodness.
\end{theorem}

\begin{proof}
The indicator of the union is bounded by the sum of indicators:
\[
 [E_1(i)\vee\cdots\vee E_m(i)]
 \leq\sum_j[E_j(i)].
\]
Sum over \(i<L\) using a bounded fold to obtain a union count at most
\(\sum_jb_j<L\).  Therefore not every index is exceptional.  Apply
Theorem~\ref{thm:least-search} to the complement predicate.
\end{proof}

This theorem combines finitely many variation, oscillation, bad-prime, or
singular-index estimates.  It is often the correct constructive replacement
for an informal diagonal subsequence argument.  The output is a finite stable
window for the requested parameters, not a global convergence rate unless the
budgets themselves provide one.

\subsection{Height bounds and complete lists}

Let \(P(z,x)\) be a decidable predicate on coded objects \(x\), and
\(h(x)\) a natural-valued height.  Suppose a primitive-recursive enumerator
\(E(B)\) lists every valid code of height at most \(B\).  For rationals, one
can enumerate reduced numerator-denominator pairs; for integer tuples,
bounded boxes; for polynomials, bounded coefficient arrays.

\begin{theorem}[Height-to-list]
\label{thm:height-to-list}
If iPRA proves
\[
 P(z,x)\longrightarrow h(x)\leq B(z),
\]
then filtering \(E(B(z))\) by the Boolean checker for \(P\) yields a
primitive-recursive complete finite list of all \(x\) satisfying \(P(z,x)\).
The list comes with embedded soundness and completeness proofs.
\end{theorem}

\begin{proof}
Every listed output passes the predicate checker, so soundness follows from
its reflection theorem.  Conversely, if \(P(z,x)\), the height theorem gives
\(h(x)\leq B(z)\); enumeration completeness places an equivalent code for
\(x\) in \(E(B(z))\), and the predicate reflection ensures that filtering
retains it.  All scans are bounded by the explicit finite enumeration size.
\end{proof}

\begin{theorem}[List-to-height]
\label{thm:list-to-height}
Suppose a finite list \(L(z)\) is proved complete for \(P(z,-)\).  Then
\[
 B_L(z)=\max\{h(x):x\in L(z)\},
\]
with default zero for the empty list, is a primitive-recursive height bound for
all solutions.
\end{theorem}

\begin{proof}
Completeness gives, for every solution \(x\), a list position containing an
equivalent code.  Height invariance under code equivalence and the defining
property of the finite maximum yield \(h(x)\leq B_L(z)\).
\end{proof}

Thus, for decidable coded sets with an explicit bounded-height enumerator, a
proved height bound and a proved complete finite list are interconvertible
interfaces.  This is an important simplification: theorem packets may choose
whichever producer is mathematically natural.

\subsection{Why cardinality is different}

A statement \(\#\{x:P(z,x)\}\leq N\) does not supply a height bound.  The
solutions can be arbitrarily sparse and delayed.  A uniform conversion from
cardinality to cutoff would decide halting in standard coding examples: make a
singleton appear at the halting time of a machine, or leave the set empty if it
never halts.  The set always has cardinality at most one, but a uniform cutoff
containing it would decide whether and when the machine halts.

\begin{barrier}[Cardinality is not a search resource]
The PCOG compiler accepts a height/cutoff theorem, a complete list, or a finite
cover whose exhaustive expansion has a bound.  It does not accept a bare
cardinality bound as sufficient input to a complete search node.
\end{barrier}

This boundary pinpoints the missing component in several arithmetic-geometry
packets.  A zero-count or dimension argument may bound the number of points in
a section; it does not by itself enumerate the rational points.

\subsection{Finite auxiliary exhaustion}

Sometimes every target object admits some auxiliary certificate in a finite
box, but there is no uniform selector choosing the auxiliary from the target.
The correct construction exhausts the auxiliary box.

Let \(A<B(z)\) range over all auxiliary codes.  For each \(A\), suppose a
complete finite list \(L_A(z)\) of objects compatible with \(A\) can be
computed.  If every target \(x\) has at least one good auxiliary
\(A<B(z)\), then the union of all \(L_A(z)\), filtered by the target predicate,
is complete.  No function \(x\mapsto A_x\) is needed.

This method applies to finite polynomial auxiliaries, local charts, modular
lifts, descent coverings, and finite families of subspaces.  The coverage
theorem---that every target has a good auxiliary in the finite box---remains
the substantive obligation.

\subsection{Certificate-producing termination}

A terminating trace should ideally export not just its output but a certificate
of why it stopped: rank zero, residual below threshold, no remaining
candidate, completed table length, or checker success.  The invariant then
proves:
\begin{enumerate}
\item the rank stays within its envelope;
\item every nonterminal step lowers rank;
\item fuel exceeds the initial rank;
\item therefore the terminal flag is set by the end;
\item the terminal certificate implies the public postcondition.
\end{enumerate}

This pattern avoids reliance on host evaluation of the fuel.  The arithmetic
proof knows that the run has reached a certified terminal state.

\subsection{Combining resource nodes}

Resource bounds compose in the same graph as precision.  A polynomial
factorization node may need:
\begin{itemize}
\item coefficient precision from a backward demand;
\item a degree bound from input validity;
\item a descent fuel from Lemma~\ref{lem:mixed-radix};
\item a root-height cutoff from a local theorem;
\item an exceptional-index window from Theorem~\ref{thm:finite-window}.
\end{itemize}
These become predecessor outputs.  The factorization trace begins only after
the graph has produced proofs of all of them.  This explicit dependency
structure is the primary benefit of PCOG over an informal ``effective'' proof.
